\documentclass[11pt]{article}
\usepackage{amsmath,amssymb,amsthm,mathtools}
\usepackage{hyperref}
\usepackage{cleveref}
\newtheorem{theorem}{Theorem}[section]
\newtheorem{lemma}[theorem]{Lemma}
\theoremstyle{definition}
\newtheorem{definition}[theorem]{Definition}
\DeclarePairedDelimiter\abs{\lvert}{\rvert}
\begin{document}
\tableofcontents
\section{Primes}\label{sec:primes}
\begin{definition}\label{def:prime}
An integer $p > 1$ is \emph{prime} if its only positive divisors are $1$ and $p$.
\end{definition}
\begin{theorem}[Euclid]\label{thm:infinite}
There are infinitely many primes.
\end{theorem}
\begin{proof}
Suppose $p_1, \dots, p_n$ are all the primes and let $N = \prod_{i=1}^{n} p_i + 1$.
Then no $p_i$ divides $N$, contradicting \cref{def:prime}.
\end{proof}
\section{Analysis}
\begin{lemma}\label{lem:bound}
For all $x \in \mathbb{R}$, $\abs*{\sin x} \le \abs{x}$.
\end{lemma}
\begin{align}
  \zeta(s) &= \sum_{n=1}^{\infty} \frac{1}{n^s} = \prod_{p \text{ prime}} \frac{1}{1 - p^{-s}}, \label{eq:euler}\\
  \int_{-\infty}^{\infty} e^{-x^2}\,dx &= \sqrt{\pi}.
\end{align}
By \cref{thm:infinite} and \eqref{eq:euler}, see also \cref{lem:bound} in \cref{sec:primes}.
\newpage
\section{Matrices}
\[
  \det\begin{pmatrix} a & b \\ c & d \end{pmatrix} = ad - bc,\qquad
  \begin{cases} x + y = 1 \\ x - y = 0 \end{cases}
\]
\end{document}
